\documentclass[a4paper,12pt]{report}
\usepackage[a4paper,left=3cm,right=3cm,tmargin = 3cm]{geometry}
%\linespread{1.6}
\usepackage{leading}
%\leading{1 cm}

\usepackage{xcolor}
\definecolor{orange}{rgb}{1,0.5,0}
\usepackage{amssymb,amsmath,alltt}
\usepackage{graphicx}
\usepackage[latin1]{inputenc}
\usepackage[spanish]{babel}
\usepackage{setspace}
\usepackage{epsfig}
\usepackage{lscape}
\usepackage{subfigure}

\usepackage{longtable}
\usepackage[T1]{fontenc}

\def\changemargin#1#2{\list{}{\topmargin#2\leftmargin#1}\item[]}
\let\endchangemargin=\endlist


\DeclareMathOperator{\supess}{supess}
%%%%%%%%%%%%%%%%%%%%%%%%%%
\usepackage{amsmath,amssymb,latexsym,amsthm,amscd}
%\usepackage[ansinew]{inpute

\usepackage{times}
%\hyphenation{Python te-ne-mos cons-tan-tes Windows pro-ble-ma re-qui-si-tos o-pe-ra-do-res e-qui-li-brio ob-te-ne-mos ob-te-ni-do di-fe-ren-cial Fa-cul-tad}
%\hyphenation{mo-de-la li-mi-ta-do ge-ne-ral-men-te pue-den es-ta-bi-li-dad des-cri-ba e-xis-ten-ce e-va-l\'u-e a-pro-xi-ma-ci-\'on}
%\hyphenation{a-sin-t\'o-ti-ca-men-te Newton ins-tan-te ge-ne-ra-dor ve-ri-fi-car u-ni-di-men-sio-nal ex-pe-ri-men-ta sa-tis-fa-cer li-ne-al de-sa-rro-llar}
%\hyphenation{mo-di-fi-car}

\usepackage{psfrag}
%\usepackage[hypertex]{hyperref}[2003/11/30]
%\usepackage{lscape}|
%\usepackage{afterpage}
\usepackage{float}
%\usepackage{longtable}
%\usepackage{marvosym}
%\usepackage{titlesec}

\usepackage{verbatim}

\usepackage[usenames,dvipsnames]{pstricks}
\usepackage{epsfig}
\usepackage{pst-grad} % For gradients
\usepackage{pst-plot} % For axes

\newtheorem{theorem}{Teorema}[section]
\newtheorem{acknowledgement}[theorem]{Agradecimiento}
\newtheorem{algorithm}[theorem]{Algoritmo}
\newtheorem{assumption}{Suposici\'on}[section]
\newtheorem{axiom}[theorem]{Axioma}
\newtheorem{case}[theorem]{Caso}
\newtheorem{claim}[theorem]{Afirmaci\'{o}n}
\newtheorem{conclusion}[theorem]{Conclusi\'{o}n}
\newtheorem{condition}[theorem]{Condici\'{o}n}
\newtheorem{conjecture}[theorem]{Conjetura}
\newtheorem{corollary}[theorem]{Corolario}
\newtheorem{criterion}[theorem]{Criterio}
\newtheorem{definition}[theorem]{Definici\'{o}n}
\newtheorem{exercise}[theorem]{Ejercicio}
\newtheorem{lemma}[theorem]{Lema}
\newtheorem{notation}[theorem]{Notaci\'{o}n}
\newtheorem{problem}[theorem]{Problema}
\newtheorem{proposition}[theorem]{Proposici\'{o}n}
\newtheorem{remark}{Oservaci\'on}
\newtheorem{solution}[theorem]{Soluci\'{o}n}
\newtheorem{summary}[theorem]{Resumen}
%\newenvironment{demos}[1][Demos]{\noindent\textbf{#1.} }

\newtheorem{def_app}{Definici\'{o}n}
\newtheorem{lem_app}{Lema}[section]
\newtheorem{teo_app}{Teorema}
\newtheorem{demos_lem}[]{Demostraci\'{o}n}
\newtheorem{demos}[]{Demostraci\'{o}n}
\newtheorem{coro}{Corolario}
\newtheorem{demos_coro}[]{Demostraci\'{o}n}

%\newenvironment{prueba}[1][Prueba]{\noindent\textbf{#1.} }{\ \rule{0.5em}{0.5em}}
%\newtheorem
\newenvironment{prueba}[1][Demostraci\'{o}n]{\noindent\textbf{#1.} }

%%%%%%%%%%%%%%%%%%%%%%%%%


% Copyright (C) 2000 Sergio Mendoza sergio@mrao.cam.ac.uk
% Astrophysics Group
% Cavendish Laboratory
% Cambridge UK
%
% This program is free software; you can redistribute it and/or modify
% it under the terms of the GNU General Public License as published by
% the Free Software Foundation; either version 2 of the License, or
% (at your option) any later version.
%
% This program is distributed in the hope that it will be useful,
% but WITHOUT ANY WARRANTY; without even the implied warranty of
% MERCHANTABILITY or FITNESS FOR A PARTICULAR PURPOSE. See the
% GNU General Public License for more details.
%
% You should have received a copy of the GNU General Public License
% along with this program; if not, write to the Free Software
% Foundation, Inc., 59 Temple Place, Suite 330, Boston, MA 02111-1307 USA

% Last modified: 21/July/2008 by Esteban Ricalde e_ricalde@yahoo.com.mx
\usepackage{fancyhdr}
\usepackage{geometry}
\usepackage[small,bf]{caption}

%\pagestyle{fancy}
%
%% Coloca una l�nea en los encabezados
%\fancyhf{}
%\fancyhead[LE,RO]{\bfseries\thepage}
%\fancyhead[LO]{\bfseries\rightmark}
%\fancyhead[RE]{\bfseries\leftmark}
%
%% Cambia la estructura de una p�gina en blanco
%\fancypagestyle{plain}{
%\fancyhead{}
%\renewcommand{\headrulewidth}{0pt}
%}

% Cambia el ancho del encabezado
\setlength{\headwidth}{16.5cm}

% Cambia el espacio para el encabezado
\setlength{\headheight}{20pt}

% Cambia el margen de los pies de figura
\setlength{\captionmargin}{20pt}

% Borra la palabra Cap�tulo del \chaptermark:
\renewcommand{\chaptermark}[1]{\markboth{\MakeUppercase{\thechapter. #1}}{}}

% Define comando para colocar p�ginas en blanco antes de iniciar cada cap�tulo
\newcommand{\clearemptydoublepage}{\newpage{\pagestyle{empty}\cleardoublepage}}
\newcommand{\HRule}{\rule{\linewidth}{0.5mm}}


%\includeonly{aprobacion,catalografica,dedicatoria,resumen,agradecimientos,3.capitulo01,3.capitulo02}

\begin{document}

\makeatletter
\def\@roman#1{\romannumeral #1}
\makeatother

\pagenumbering{roman}\setcounter{page}{2}    

\begin{titlepage}
\setlength{\parindent}{0pt} \setlength{\parskip}{0pt}
%\pagenumbering{roman} \setcounter{page}{i}
\begin{center}

\textsc{\large \textbf{UNIVERSIDAD NACIONAL MAYOR DE SAN MARCOS}}\\[0.2cm]
{\large \textbf{FACULTAD DE CIENCIAS MATEM�TICAS}}\\[2cm]
\end{center}

\begin{center}

\includegraphics[scale = 0.1]{UNMSM_logo}
\end{center}

\begin{center}

\vfill
{\Large \bfseries  Modelamiento de un sistema termoel�stico unidimensional y su simulaci�n computacional usando Python}\\[0.4cm]
%{\Large \bfseries  Modelado y simulaci�n computacional de un sistema termoel�stico usando python c�mo herramienta de simulaci�n}\\[0.4cm]
% {\Large \bfseries Deducci�n y estudio de un sistema termoel�stico\\ e implementaci�n en base a software libre}\\[0.4cm]
\end{center}

\vspace*{\stretch{1}}% \HRule
\begin{center}
Tesis para Optar el T�tulo Profesional de \\[0.4cm]
Licenciado en Computaci�n Cient�fica \\[2cm]
\end{center}



\begin{center}
\textbf{LUIS ALONSO MANSILLA ALVAREZ}\\[0.4cm]
\end{center}
%\HRule 
\vspace*{\stretch{2}}

\begin{center}
LIMA - PER� \\ Diciembre - 2011
\end{center}
%vfill{}
\end{titlepage}
\clearemptydoublepage

\vspace{3cm}
\begin{center}

%\vfill
%\vspace*{\stretch{1}}% \HRule
\underline{MODELAMIENTO DE UN SISTEMA TERMOEL�STICO UNIDIMENSIONAL Y}\\ \underline{SU SIMULACI�N COMPUTACIONAL USANDO PYTHON}\\[1.5cm]
%\underline{DEDUCCI�N Y ESTUDIO DE UN SISTEMA TERMOEL�STICO}\\ \underline{E IMPLEMENTACI�N EN BASE A SOFTWARE LIBRE}\\[1.5cm]
\end{center}

\begin{center}
Luis Alonso Mansilla Alvarez \\[1.5cm]
\end{center}

Tesis presentada a consideraci�n del Cuerpo Docente de la Facultad de Ciencias Matem�ticas, de la Universidad Nacional Mayor 
de San Marcos, como parte de los requisitos para obtener el T�tulo Profesional de Licenciado en Computaci�n Cient�fica.\\

Aprobada por:\\
\vspace{1cm}
\begin{center}
.........................................................\\
Dr. Manuel Efra�n Carbajal Pe�a\\
\quad Presidente \\[1.5cm]
%Dr. Manuel Efra�n Carbajal Pe�a
..........................................................\\
Dra. Roxana L�pez Cruz\\
\quad Miembro \\[1.5cm]
%El otro px
...........................................................\\
Dr. Jos� Ra�l Luyo S�nchez\\
\quad Miembro Asesor \\[1.5cm]
%DSc. Jos� Raul Luyo S�nchez
\end{center}

\vfill
\begin{center}
	Lima - Per� \\ Diciembre - 2011
\end{center}

\clearemptydoublepage
\newenvironment{dedication} {\cleardoublepage
%\thispagestyle{empty}
\vspace*{\stretch{1}} \begin{center}} {\end{center}
\vspace*{\stretch{3}} \clearpage}
\begin{dedication}
\begin{center}
FICHA CATALOGR�FICA\\[4cm]
\end{center}
\hfill \
\parbox{10cm}{
\hspace{-1cm}MANSILLA ALVAREZ, LUIS ALONSO\\[0.1cm]
Modelamiento de un sistema termoel�stico unidimensional y su simulaci�n computacional usando Python, (Lima) 2011.\\
viii, 82 p., 29,7 cm, (UNMSM, Licenciado, Computaci�n Cient�fica, 2011).\\
Tesis, Universidad Nacional Mayor de San Marcos, Facultad de Ciencias Matem�ticas 1. Computaci�n Cient�fica. UNMSM/FdeCM II. T�tulo (Serie).\\[8cm]

}
\end{dedication}
%\clearemptydoublepage
\thispagestyle{empty}
\newenvironment{dedicationo} {\cleardoublepage
%\thispagestyle{empty}
\vspace*{\stretch{1}} \begin{flushright} \em} {\end{flushright}
\vspace*{\stretch{3}} \clearpage}

\begin{dedicationo}
\begin{minipage}[c]{6cm} 
\,
\end{minipage}
\begin{minipage}[c]{8cm}
 Dedicatoria
\end{minipage}
\end{dedicationo}
%\clearemptydoublepage
\chapter*{Agradecimientos
\markboth{Agradecimientos}{Agradecimientos}}
Agradecimientos
\vfill
\hfill {\it Luis Alonso Mansilla Alvarez}

%\leading{0.8 cm}
\clearemptydoublepage
%%\pagenumbering{roman} \setcounter{page}{1}
%\chapter*{Resumen\markboth{Resumen}{Resumen}}
%Vendan bien sus tesis
\begin{center}
RESUMEN\\[0.7cm]
MODELAMIENTO DE UN SISTEMA TERMOEL�STICO UNIDIMENSIONAL Y SU SIMULACI�N COMPUTACIONAL USANDO PYTHON\\[0.7cm]
%DEDUCCI�N Y ESTUDIO DE UN SISTEMA TERMOEL�STICO\\ E IMPLEMENTACI�N EN BASE A SOFTWARE LIBRE\\[0.7cm]
LUIS ALONSO MANSILLA ALVAREZ\\[0.7cm]
DICIEMBRE - 2011\\[0.7cm]
\end{center}

\hspace{-0.6cm}Orientador: \qquad\qquad\qquad DSc. Jos� Raul Luyo S�nchez\\
T�tulo obtenido: \qquad\qquad\,Licenciado en Computaci�n Cient�fica\\
\noindent{.\dotfill{}.\par}
%--------------------------------------------------------------------------\\
\hspace{-0.6cm}En este trabajo presentamos un estudio completo de la termoelasticidad lineal en el caso unidimensional. Basados en leyes f�sicas, 
desarrollamos un modelo matem�tico que describa dicho fen�meno en base a un sistema de ecuaciones diferenciales parciales que involucren el desplazamiento
espacial y la temperatura.\\
\hspace{-0.6cm}Planteado el modelo matem�tico, obtenemos los resultados que garantizan que el problema est� ``bien puesto'' en el sentido que propuso 
Hadamard, aseguramos para esto la existencia, unicidad y regularidad de la soluci�n en base a la teor�a de semigrupos de operadores aplicada a las 
ecuaciones diferenciales parciales.\\
\hspace{-0.6cm}Como punto central, desarrollamos un esquema num�rico para simular el comportamiento del fen�meno y proponemos el uso del lenguaje de 
programaci�n Python como tecnolog�a de software libre, �til y efectiva para la implementaci�n computacional y obtenci�n de resultados mediante el esquema 
num�rico. \\[1cm]
\begin{tabular}{ l l }
  \hspace{-0.2cm}PALABRAS CLAVES: & ECUACIONES DIFERENCIALES PARCIALES\\
  & TERMOELASTICIDAD \\
  & SEMIGRUPOS \\
  & PYTHON \\
  & ALGORITMOS\\
  
\end{tabular}
\clearemptydoublepage
%%\pagenumbering{roman} \setcounter{page}{1}
%\chapter*{Resumen\markboth{Resumen}{Resumen}}
%Vendan bien sus tesis
\begin{center}
ABSTRACT\\[1cm]
MODELING OF A LINEAR THERMOELASTIC SYSTEM AND ITS COMPUTATIONAL SIMULATION USING PYTHON\\[0.7cm]
%DEDUCTION AND STUDY OF A THERMOELASTIC SYSTEM AND\\ HIS IMPLEMENTATION BASED ON FREE SOFTWARE\\[0.7cm]
LUIS ALONSO MANSILLA ALVAREZ\\[0.7cm]
DECEMBER - 2011\\[0.7cm]
\end{center}

\hspace{-0.6cm}Advisor: \qquad\qquad\qquad\,\,\, DSc. Jos� Raul Luyo S�nchez\\
Obtained Title: \qquad\qquad\,Degree in Scientific Computing\\
\noindent{.\dotfill{}.\par}
\hspace{-0.6cm}We present a complete study of the linear thermoelasticity in one-dimensional case. Based on physical laws, we develope a mathematical model 
which describes this phenomenon based on a system of partial differential equations involving the spatial displacement and temperature.\\
\hspace{-0.6cm}Once retrieved the mathematical model, we obtain the results that guarantee that it is a ``well-posed'' problem in the sense proposed by 
Hadamard, for this, we ensure the existence, uniqueness and regularity of the solution based on the theory of semigroups of operators applied to partial 
differential equations.\\
\hspace{-0.6cm}As central point, we develop a numerical scheme to simulate the behaivor of the phenomenon and propose to use the programming language Python as 
an useful and effective free software for computational implementation and achivement of results by the numerical scheme.\\[1cm]
\begin{tabular}{ l l }
  \hspace{-0.2cm}KEYWORDS: & PARTIAL DIFFERENTIAL EQUATIONS\\
  & THERMOELASTICITY \\
  & SEMIGROUPS \\
  & PYTHON \\
  & ALGORITHMS\\
  
\end{tabular}

%\leading{1 cm}
\clearemptydoublepage
\tableofcontents

\chapter{PRODUCTO VECTORIAL}
%{\ref{producto.vectorial} \sc Producto vectorial}

\section{ Producto vectorial.}
\label{producto.vectorial}

\bigskip
Consideremos la terna ordenada de vectores de $V$, $\{
\overrightarrow{i} ,\overrightarrow{j} ,\overrightarrow{k} \}$,
que constituye una base de $V$. Estas las vamos a clasificar en
dos grupos: uno formado por las ternas $\left\{ {i^{} ,j,k }
\right\}$ tales que para un observador parado en $v_3$, la
rotaci\'{o}n de \'{a}ngulo convexo (o sea, de medida $< \pi $) que
lleva $i$ en $j$ es de sentido antihorario; el otro formado por
las ternas en las que esa rotaci\'{o}n tiene sentido horario.
% Warning: EPS-printer was not specified  -- Figure omitted!


\begin{figure}[htb]
\centering
FIGURA
%\includegraphics[scale=0.4]{orienta.eps}
\end{figure}

\bigskip

Para definir el producto vectorial necesitamos elegir uno de los
dos tipos como terna preferida. Nos quedaremos para esto con las
ternas del primer tipo que llamaremos \textbf {positivas}.
\bigskip

Esta convenci\'{o}n sirve para definir el producto vectorial como
una funci\'{o}n de $V \times V\,\,en\,\,V$ (as\'{\i} como el
producto escalar era una funci\'{o}n de $V \times
V\,\,\,en\,\,\,\,\mathbb{R}$) que a cada par de vectores $v,w$ le
asocia un vector, que denotaremos $v$ $ \wedge $ $w$, que cumple:

\noindent a) $\left\| {v \wedge w} \right\| = \left\| {v}
\right\|.\left\| {w} \right\|.sen\,\theta $  ($\theta $ \'{a}ngulo
de \textit{v} \,\, con \,\, \textit{w} )

\noindent b) $\left( {v \wedge w} \right).v = 0\,\,y\,\,\left( {v
\wedge w} \right).w = 0$

\noindent c) Si $v \ne 0\,\,\,\,y\,\,\,\,w \ne 0$ la terna
$\left\{ {v,w,v \wedge w} \right\}$ es positiva.


\bigskip

\textbf{Propiedades:}

\textbf{1)} $\left\| {v \wedge w} \right\|$ es el doble del
\'{a}rea del tri\'{a}ngulo determinado por esos vectores.


\textbf{2)} $v \wedge w = 0$ si y s\'{o}lo si \textit{v }y w son
colineales (en particular, si alguno de ellos es el vector $0$).


\textbf{3) }Respecto de la condici\'{o}n  c) de la definici\'{o}n,
corresponde observar que si $v \wedge w \ne 0$ la terna $\left\{
{v,w,v \wedge w} \right\}$ es una base, pues por b) esos vectores
no son coplanares salvo que \textit{v }y w sean colineales, en
cuyo caso $v \wedge w = 0$

\textbf{4) }$v \wedge w =  - \left( {w \wedge v} \right)$ ( en
particular , esta operaci\'{o}n no es conmutativa) Para verificar
esto basta notar que si para cierto observador la rotaci\'{o}n del
\'{a}ngulo $ < \pi $ que lleva \textit{v }en w es de sentido
antihorario, para el mismo observador la que lleva w en \textit{v
}es de sentido horario. De modo que el observador debe ubicarse
del lado opuesto del plano u,w para que esa rotaci\'{o}n aparezca
de sentido trigonom\'{e}trico . Luego esa debe ser la
ubicaci\'{o}n de $w \wedge v$ para que la terna $\left[ {w,v,w
\wedge v} \right]$ sea positiva.

\begin{figure}[htb]
\centering
FIGURA
%\includegraphics[scale=0.5]{vectorial.eps}
\end{figure}
%\begin{center}
%% Warning: EPS-printer was not specified  -- Figure omitted!
%
%\end{center}
\textbf{5) }Este producto no es asociativo. Se puede deducir que:
\[
\left( {u \wedge v} \right) \wedge w + \left( {v \wedge w} \right)
\wedge u + \left( {w \wedge u} \right) \wedge v = 0
\] de donde, usando 4), $\left( {u \wedge v} \right) \wedge w = u
\wedge \left( {v \wedge w} \right) + v \wedge \left( {w \wedge u}
\right)$, y como en general $v \wedge \left( {w \wedge u} \right)
\ne 0$, la propiedad asociativa no vale.

\textbf{6) }Para todo $\lambda  \in \mathbb{R}$: $\lambda \left(
{v \wedge w} \right) = \left( {\lambda v} \right) \wedge w = v
\wedge \left( {\lambda w} \right)$. Esta propiedad se deduce
directamente de la definici\'{o}n en el caso $\lambda  \ge
0$. Para el caso $\lambda  < 0$, hay que observar que del hecho
que la terna $\left\{ {v,w,v \wedge w} \right\}$ es positiva se
deduce que $\left\{ { - v,w, - \left( {v \wedge w} \right)}
\right\} $ y
\\ $ \left\{ {v, - w, - \left( {v \wedge w} \right)} \right\}$ son
tambi\'{e}n positivas.

\textbf{7) }El producto vectorial es distributivo respecto de la
suma de vectores. Esto es:

\noindent a) $\left( {v_1  + v_2 } \right) \wedge w = v_1 \wedge w
+ v_{2}^{}  \wedge w.$


\noindent b) $v \wedge \left( {w_{1}^{}  + w_{2}^{} } \right) = v
\wedge w_{1}^{}  + v \wedge w_{2}^{}. $

\noindent Para demostrar esta propiedad hacemos algunas
observaciones previas. En primer lugar, observamos que llamando
$\pi $ al plano con vector normal \textit{v }e indicando por $p $
\textit{w} la proyecci\'{o}n de \textit{w} sobre ese plano, y con
$r(p\textit{w})$ el vector que se obtiene rotando esa
proyecci\'{o}n un \'{a}ngulo de medida $\pi /2$ en sentido
antihorario observado desde $v$, se verifica que: $v \wedge w =
\left\| {v} \right\|.r\left( {pw} \right)$ pues
 $\left\| {r\left( {pw} \right)} \right\| = \left\| {pw} \right\| = \left\|
{w} \right\|.sen\,\theta $ (ver figura)

\begin{figure}[htb]
\centering
FIGURA
%\includegraphics[scale=0.4]{propvec.eps}
\end{figure}
%\begin{center}
%% Warning: EPS-printer was not specified  -- Figure omitted!
%
%\end{center}


\bigskip

%\begin{center}
%FIGURA 4.1
%\end{center}


\bigskip

Como $p\left( {w_{1}^{}  + w_{2}^{} } \right) = pw_{1}^{}  +
pw_{2}^{} $ y $r\left( {pw_{1}^{}  + pw_{2}^{} } \right) = r\left(
{pw_{1}^{} } \right) + r\left( {pw_{2}^{} } \right)$ tendremos
que: $v \wedge \left( {w_{1}^{} \wedge w_{2}^{} } \right) =
\left\| {v} \right\|.r\left( {p\left( {w_{1}^{} + w_{2}^{} }
\right)} \right)$$ = \left\| {v} \right\|.r\left( {pw_{1}^{} }
\right) + \left\| {v} \right\|.r\left( {pw_{2}^{} } \right) = v
\wedge w_{1}^{}  + v \wedge w_{2}^{} $.


As\'{\i} se prueba b) y an\'{a}logamente se obtiene a).

\textbf{8) }Si $ \{i,j,k\}$ es una base ortonormal positiva de $V$
, entonces:
\[
{\begin{array}{l}
 i \wedge i = j \wedge j = k \wedge k = 0; \\
 i \wedge j = k;\;\;j \wedge k = i;\;\;k \wedge i = j \\
  \\
 \end{array}}
\]
(es decir, en la sucesi\'{o}n ($i,j,k,i,j, \cdots$) el producto de
dos vectores sucesivos da el que le sigue):
\[
j \wedge i =  - k,\,k \wedge j =  - i;\,i \wedge k =  - j
\]

\begin{figure}[htb]
\centering
FIGURA
%\includegraphics[scale=0.5]{ijk.eps}
\end{figure}

\textbf{9) }De 4.6), 4.7) y 4.8) resulta que si $v = a_{1}^{} i +
a_{2}^{} j + a_{3}^{} k$ y $w = b_{1}^{} i + b_{2}^{} j + b_{3}^{}
k$.
 entonces:
\[
v \wedge w = \left( {a_{2}^{} b_{3}^{}  - a_{3}^{} b_{2}^{} }
\right)i - \left( {a_{1}^{} b_{3}^{}  - a_{3}^{} b_{1}^{} }
\right)j + \left( {a_{1}^{} b_{2}^{}  - a_{2}^{} b_{1}^{} }
\right)k.
\]
Obs\'{e}rvese que este resultado puede recordarse pensando en el
desarrollo por la primer fila de un determinante:
\[
\left| {{\begin{array}{*{20}c}
  {i} \hfill & {j} \hfill & {k} \hfill  \\
  {a_{1}^{} } \hfill & {a_{2}^{} } \hfill & {a_{3}^{} } \hfill  \\
  {b_{1}^{} } \hfill & {b_{2}^{} } \hfill & {b_{3}^{} } \hfill  \\
\end{array} }} \right| = \left( {a_{2}^{} b_{3}^{}  - a_{3}^{} b_{2}^{} }
\right)i - \left( {a_{1}^{} b_{3}^{}  - a_{3}^{} b_{1}^{} }
\right)j + \left( {a_{1}^{} b_{2}^{}  - a_{2}^{} b_{1}^{} }
\right)k
\]
Esta expresi\'{o}n del producto vectorial puede tambi\'{e}n
tomarse como su definici\'{o}n. M\'{a}s precisamente, dada una
terna ortonormal cualquiera \{\textit{i,j,k}\} (positiva o
negativa), puede definirse el producto de dos vectores por la
expresi\'{o}n dada en 4.9). Se comprueba entonces sin dificultad
que el vector:
 $u = \left( {a_{2}^{} b_{3}^{}  - a_{3}^{} b_{2}^{} } \right)i - \left(
{a_{1}^{} b_{3}^{}  - a_{3}^{} b_{1}^{} } \right)j + \left(
{a_{1}^{} b_{2}^{}  - a_{2}^{} b_{1}^{} } \right)k$ verifica
$\left\| {u} \right\| = \left\| {v} \right\|.\left\| {w}
\right\|.sen\,\theta $ y $u.v = u.w = 0$. Si adem\'{a}s $
\{i,j,k\}$ es una terna positiva puede verificarse, que $
\{v,w,u\}$ es tambi\'{e}n positiva. Luego $u = v \wedge w$ , por
lo tanto esta definici\'{o}n de $v \wedge w$ coincide con la
inicial.
\bigskip

\subsection{Aplicaciones geom\'{e}tricas.}

\textbf{a ) }\textbf {Distancia de un punto a una recta:} sea r
una recta dada por un punto A y un vector \textit{v}. La distancia
de Q a r es:
\[
d\left( {Q,r} \right) = \left| {d\left( {Q,A} \right).sen\,\theta
} \right| = \left\| {AQ} \right\|.sen\,\theta  = \frac{{\left\|
{AQ \wedge v} \right\|}}{{\left\| {v} \right\|}}
\]


\begin{figure}[htb]
\centering
FIGURA
%\includegraphics[scale=0.4]{puntorecta.eps}
\end{figure}
%\begin{center}
%% Warning: EPS-printer was not specified  -- Figure omitted!
%
%\end{center}
%
%
%\bigskip
%
%\begin{center}
%FIGURA 5.1
%\end{center}


\bigskip

Si $A = ( x_{o} ,y_{o} ,z_{o});\,\,Q = (\overline {x} ,\overline
{y} ,\overline {z} );\,v = ai\, + \,bj\, + \,ck$ entonces la
distancia es $ d(Q,r)=$ $$
\frac{\sqrt{[(\overline{y}-y_{0})c-(\overline{z}-z_{0})b]^{2}+
[(\overline{x}-x_{0})c-(\overline{z}-z_{0})a]^{2}+[(\overline{x}-x_{0})b-(\overline{y}-y_{0})a]^{2}}}
{\sqrt{a^{2}+b^{2}+c^{2}}}$$

\bigskip

\begin{ejemplo}\label{piej5.1}En el punto 3 de este cap\'{\i}tulo se vieron
algunas superficies de revoluci\'{o}n particulares. Sea $r\left\{
\begin{array}{l}
 {x = y} \\
 {z = 0}
 \end{array} \right.$ $C:\left\{ \begin{array}{l}
 {x = 0} \\
 {zy = 1}
 \end{array} \right.$. Hallar la superficie de revoluci\'{o}n de eje r
y generatriz C.

Las ecuaciones son:
\[
\left\{ \begin{array}{l}
 {x_{0}^{}  = 0                                 ,\,\,\,            P_{0}^{}  \in C}
\\
 {z_{0}^{} y_{0}^{}  = 1                  ,\,\,\,                        P_{0}^{}
\in C} \\
 {x - x_{0}^{}  + y - y_{0}^{}  = 0              ,\,\,\,            plano ,\,\,\, \pi ,\,\,\, por
,\,\,\, P_{0}^{}  \bot r} \\
 {2z^{2} + \left( {x - y} \right)^{2} = 2z_{0}^{2}  + \left( {x_{0}^{}  -
y_{0}^{} } \right)^{2} ,\,\,\,  dist\left( {P,r} \right) =
dist\left( {P_{0}^{} ,r} \right)}
 \end{array} \right.
\]
eliminando $x_{0}^{} ,y_{0}^{} ,\,z_{0}^{} $ resulta $\left( {x +
y} \right)^{2}\left( {z^{2} - 2xy} \right) = 1$
\end{ejemplo}

\vspace{.5cm}

b)\textbf{Intersecci\'{o}n de dos planos no paralelos:}

Sea $r:\left\{ \begin{array}{l}
 {ax + by + cz + d = 0} \\
 {a'x + b'y + c'z + d' = 0}
 \end{array} \right.$ una recta (dada como intersecci\'{o}n de dos
planos no paralelos). Supongamos que queremos las ecuaciones
param\'{e}tricas de esa recta, es decir, las de la forma: $x =
x_{0}^{}  + \lambda p,\,\,y = y_{0}^{}  + \lambda q,\,\,z =
z_{0}^{}  + \lambda r$ (o tambi\'{e}n $P = P_{0}^{}  + \lambda
v$). Para esto se necesita hallar un punto $P_{0}^{}  = \left(
{x_{0}^{} ,y_{0}^{} ,z_{0}^{} } \right)$ de la recta y un vector
$v = p\,i + q\,j + r\,k$ de la direcci\'{o}n de r.
\bigskip
Esto se puede hacer sin usar el producto vectorial, resolviendo el
sistema de ecuaciones
\[
\left\{ \begin{array}{l}
 {ax + by + cz + d = 0} \\
 {a'x + b'y + c'z + d' = 0}
 \end{array} \right.
\]
Usando el producto vectorial: tomar $P_{0}^{}  = \left( {x_{0}^{}
,y_{0}^{} ,z_{0}^{} } \right)$ una soluci\'{o}n particular del
sistema de ecuaciones. Si el sistema de coordenadas es ortogonal,
entonces:
 $n = ai + bj + ck\,\,y\,\,n' = a'i + b'j + c'k$ son vectores normales a los
planos, y entonces $n \wedge n'$ es un vector de la
intersecci\'{o}n de ambos planos (luego, est\'{a} en la
direcci\'{o}n de la recta de intersecci\'{o}n de ambos).


\begin{figure}[htb]
\centering
FIGURA
%\includegraphics[scale=0.4]{planos.eps}
\end{figure}
%\begin{center}
%% Warning: EPS-printer was not specified  -- Figure omitted!
%
%\end{center}
%
%\bigskip
%\begin{center}
%FIGURA 5.2
%\end{center}

Luego: $n \wedge n'$$ = \left( {bc' - b'c} \right)i - \left( {ac'
- \,ca'} \right)j + \left( {ab' - ba'} \right)k$ es de la
direcci\'{o}n de r.

\bigskip
\textbf{c ) }\textbf {Distancia entre dos rectas:} Se define la
distancia $d\left( {r,r'} \right)$ entre dos rectas $r,r'$ como el
m\'{\i}nimo de $d\left( {P,P'} \right)$ donde $P \in r$ y $P' \in
r'$.

Es claro que este m\'{\i}nimo se obtiene cuando la recta $PP'$ es
perpendicular al mismo tiempo a $r $ y a $r'$. Sea $ r $ dada por
un punto $A \in r$ y un vector $v$  de su direcci\'{o}n y $r'$ por
$B \in r$ y un vector $w$.


Para tener $d\left( {r,r'} \right)$: si \textit{r} y \textit{r'}
son paralelas: basta tomar el punto A y hallar $dist\left( {A,r'}
\right)$; si $r$ y $r'$ se cortan: $d\left( {r,r'} \right)$= 0.


\begin{figure}[htb]
\centering
FIGURA%
FIGURA%\includegraphics[scale=0.5]{rectarecta.eps}
\end{figure}
%\begin{center}
%% Warning: EPS-printer was not specified  -- Figure omitted!
%
%\end{center}
%\bigskip
%\begin{center}
%FIGURA 5.3
%\end{center}

\bigskip

Supongamos ahora que \textit{r} y \textit{r'} no son coplanares.
Si \textit{s} es la perpendicular com\'{u}n , $P_{o} \,\, =
\,\,s\,\, \cap \,\,r^{}$  y  $P_{0}^{'}  = s \cap r'$, entonces:
 $d\left( {r,r'} \right) = d\left( {P_{0}^{} ,P_{0}^{'} } \right) = d\left(
{A,\pi } \right)$ donde $\pi $  es el plano paralelo a \textit{r}
que contiene a \textit{r'}. (ver figura) . Como versor de la
direcci\'{o}n de \textit{s} puede tomarse $n = \frac{{v \wedge
w}}{{\left\| {v \wedge w} \right\|}}$. Luego :
\[
d\left( {r,r'} \right) = d\left( {A,\pi } \right) =
\frac{{1}}{{\left\| {v \wedge w} \right\|}}\left| {\left( {v
\wedge w} \right).\overrightarrow{AB} } \right|
\]


\begin{figure}[htb]
\centering
FIGURA
%\includegraphics[scale=0.5]{normalcomun.eps}
\end{figure}
%\begin{center}
%% Warning: EPS-printer was not specified  -- Figure omitted!
%
%\end{center}
%
%\bigskip
%\begin{center}
%FIGURA 5.3
%\end{center}
\bigskip

\textbf{d ) }\textbf {Perpendicular com\'{u}n a dos rectas:} Si
las dos rectas son paralelas, una perpendicular com\'{u}n es la
recta perpendicular a una de ellas trazada por un punto de la
otra. Si se interceptan , el problema es tambi\'{e}n f\'{a}cil de
resolver. Supongamos ahora que las dos rectas no son coplanares.
La perpendicular com\'{u}n $s$ est\'{a} en el plano $\pi $ de $s$
y $r$, y en el $\pi $' determinado por $s$ y $r'$. Luego $s = \pi
\cap \pi '$. Dada la direcci\'{o}n \textit{v }de \textit{r} y la
$w$ de \textit{r'}, y $A \in r,\,\,B \in r'$, el plano $\pi $
queda determinado por $A$,\textit{v} y $v \wedge w$. El $\pi $'
por B,\textit{w},$v \wedge w$, pues $v \wedge w$ es un vector de
la direcci\'{o}n de \textit{s}. Las ecuaciones de $\pi $ y $\pi $'
as\'{\i} obtenidas constituyen un par de ecuaciones que
representan la recta \textit{s}.
\bigskip

\textbf{e ) }\textbf{Volumen de un tetraedro: }Consideremos un
tetraedro dado por sus v\'{e}rtices $A, B, C, D$. Su volumen es $V
= \frac{{1}}{{3}}$ \'{a}rea de la base $ \times $ altura.
Tendremos: \'{a}rea base = $1/2\left\| {\left( {B - A} \right)
\wedge \left( {C - A} \right)} \right\|$.

\begin{figure}[htb]
\centering
FIGURA
%\includegraphics[scale=0.4]{tetra.eps}
\end{figure}

%% Warning: EPS-printer was not specified  -- Figure omitted!
%
%\bigskip
%\begin{center}
%FIGURA 5.4
%\end{center}

 Si \textit{n} es el versor normal a la base, la altura es $
 \left| {n.\left( {D - A} \right)} \right|$\, con \,$n = \frac{{\left(
{B - A} \right) \wedge \left( {C - A} \right)}}{{\left\| {\left(
{B - A} \right) \wedge \left( {C - A} \right)} \right\|}}.$ Luego
$V = 1/6.\left\| {\left( {B - A} \right) \wedge \left( {C - A}
\right)} \right\|.\frac{{\left| {\left[ {\left( {B - A} \right)
\wedge \left( {C - A} \right)} \right].\left( {D - A} \right)}
\right|}}{{\left\| {\left( {B - A} \right) \wedge \left( {C - A}
\right)} \right\|}}$
\[ {\rm por\,\,\, lo\,\,\, que\,\,\,\,}
V = 1/6.\left| {\left[ {\left( {B - A} \right) \wedge \left( {C -
A} \right)} \right].\left( {D - A} \right)} \right|
\]

\bigskip

\subsection{Producto mixto}
\bigskip
Es una operaci\'{o}n definida de $V\, \times \,\,V\,\, \times
\,\,V\,$ en $\mathbb{R}$. Dados tres vectores $u,v, \,y\,\,w$,
llamamos producto mixto al n\'{u}mero $\left( {v \wedge w}
\right).u$ que indicamos $v \wedge w.u$, o tambi\'{e}n $u.\left(
{v \wedge w} \right)$.
\bigskip
Consideremos un sistema ortonormal $\left\{ {i,j,k} \right\}$ y
sean:
\[
v = a_{1}^{} i + a_{2}^{} j + a_{3}^{} k, w = b_{1}^{} i +
b_{2}^{} j + b_{3}^{} k, u = c_{1}^{} i + c_{2}^{} j + c_{3}^{} k.
\]
Entonces:
\[
v \wedge w = \left| {{\begin{array}{*{20}c}
  {a_{2}^{} } \hfill & {a_{3}^{} } \hfill  \\
  {b_{2}^{} } \hfill & {b_{3}^{} } \hfill  \\
\end{array} }} \right|i - \left| {{\begin{array}{*{20}c}
  {a_{1}^{} } \hfill & {a_{3}^{} } \hfill  \\
  {b_{1}^{} } \hfill & {b_{3}^{} } \hfill  \\
\end{array} }} \right|j + \left| {{\begin{array}{*{20}c}
  {a_{1}^{} } \hfill & {a_{2}^{} } \hfill  \\
  {b_{1}^{} } \hfill & {b_{2}^{} } \hfill  \\
\end{array} }} \right|k.
\]
\[
v \wedge w.u = \left| {{\begin{array}{*{20}c}
  {a_{2}^{} } \hfill & {a_{3}^{} } \hfill  \\
  {b_{2}^{} } \hfill & {b_{3}^{} } \hfill  \\
\end{array} }} \right|c_{1}^{}  - \left| {{\begin{array}{*{20}c}
  {a_{1}^{} } \hfill & {a_{3}^{} } \hfill  \\
  {b_{1}^{} } \hfill & {b_{3}^{} } \hfill  \\
\end{array} }} \right|c_{2}^{}  + \left| {{\begin{array}{*{20}c}
  {a_{1}^{} } \hfill & {a_{2}^{} } \hfill  \\
  {b_{1}^{} } \hfill & {b_{2}^{} } \hfill  \\
\end{array} }} \right|c_{3}^{} =
\left| {{\begin{array}{*{20}c}
  {c_{1}^{} } \hfill & {c_{2}^{} } \hfill & {c_{3}^{} } \hfill  \\
  {a_{1}^{} } \hfill & {a_{2}^{} } \hfill & {a_{3}^{} } \hfill  \\
  {b_{1}^{} } \hfill & {b_{2}^{} } \hfill & {b_{3}^{} } \hfill  \\
\end{array} }} \right|
\]
\begin{center}
(Por desarrollo por la primer fila de ese determinante)
\end{center}
Usando el hecho de que si se permutan dos filas de una matriz
entre s\'{\i} el determinante s\'{o}lo cambia de signo, resulta
que dos de esas permutaciones no cambian el determinante; luego:
\[
\left| {{\begin{array}{*{20}c}
  {c_{1}^{} } \hfill & {c_{2}^{} } \hfill & {c_{3}^{} } \hfill  \\
  {a_{1}^{} } \hfill & {a_{2}^{} } \hfill & {a_{3}^{} } \hfill  \\
  {b_{1}^{} } \hfill & {b_{2}^{} } \hfill & {b_{3}^{} } \hfill  \\
\end{array} }} \right| = \left| {{\begin{array}{*{20}c}
  {b_{1}^{} } \hfill & {b_{2}^{} } \hfill & {b_{3}^{} } \hfill  \\
  {c_{1}^{} } \hfill & {c_{2}^{} } \hfill & {c_{3}^{} } \hfill  \\
  {a_{1}^{} } \hfill & {a_{2}^{} } \hfill & {a_{3}^{} } \hfill  \\
\end{array} }} \right| = \left| {{\begin{array}{*{20}c}
  {a_{1}^{} } \hfill & {a_{2}^{} } \hfill & {a_{3}^{} } \hfill  \\
  {b_{1}^{} } \hfill & {b_{2}^{} } \hfill & {b_{3}^{} } \hfill  \\
  {c_{1}^{} } \hfill & {c_{2}^{} } \hfill & {c_{3}^{} } \hfill  \\
\end{array} }} \right|
\]
En consecuencia: $\left( {v \wedge w} \right).u = \left( {u \wedge
v} \right).w = \left( {w \wedge u} \right).v$. Es decir, que en la
sucesi\'{o}n (\textit{u,v,w,u,v,w}) el producto vectorial de dos
vectores sucesivos multiplicado escalarmente por el siguiente, da
lo mismo cualesquiera sean los tres vectores sucesivos. Se puede
entonces hablar del producto mixto de \textit{u, v, w} sin indicar
si se trata de $\left( {u \wedge v} \right).w$ o de $u.\left( {v
\wedge w} \right)$, pues el resultado es el mismo. Es por esto que
se escribe (\textit{u,v,w}) como notaci\'{o}n para $\left( {u
\wedge v} \right).w = u.\left( {v \wedge w} \right)$.

Observamos que (\textit{v,w,u}) = 0 si y s\'{o}lo si
\'{a}ng$\left( {v \wedge w,u} \right) = \pi /2$ o alguno de los
vectores es el nulo. Como \'{a}ng $\left( {v \wedge w,v} \right) =
\,\,\mbox{\'{a}ng}\,\,\left( {v \wedge w,w} \right) = \pi /2$ ,
para que (\textit{v,w,u}) = 0 es necesario y suficiente que
\textit{v, w} y \textit{u} sean coplanares. (Obs\'{e}rvese que
esto podr\'{\i}a sacarse como conclusi\'{o}n del c\'{a}lculo del
volumen del tetraedro. $(v,w,u) = 0$ si y s\'{o}lo si el volumen
del tetraedro determinado por esos vectores es 0, esto es
equivalente a decir que los vectores son coplanares).
\bigskip

Esta condici\'{o}n permite escribir la ecuaci\'{o}n vectorial del
plano dado por tres puntos $A, B, C$ en otra forma. Decir que $P$
pertenece a ese plano equivale a decir que los vectores $P-A, B-A$
y $C-A$ son coplanares, o sea $(P-A,B-A,C-A) = 0$.


Pasando a coordenadas, si \textit{P=}(x,y,z),\textit{A=}($a_{1}^{}
,a_{2}^{} ,a_{3}^{} $), $B = \left( {b_{1}^{} ,b_{2}^{} ,b_{3}^{}
} \right)$ y $C = \left( {c_{1}^{} ,c_{2}^{} ,c_{3}^{} } \right)$
esta ecuaci\'{o}n se escribe:
\[
\left| {{\begin{array}{*{20}c}
  {x - a_{1}^{} } \hfill & {y - a_{2}^{} } \hfill & {z - a_{3}^{} } \hfill
\\
  {b_{1}^{}  - a_{1}^{} } \hfill & {b_{2}^{}  - a_{2}^{} } \hfill &
{b_{3}^{}  - a_{3}^{} } \hfill  \\
  {c_{1}^{}  - a_{1}^{} } \hfill & {c_{2}^{}  - a_{2}^{} } \hfill &
{c_{3}^{}  - a_{3}^{} } \hfill  \\
\end{array} }} \right| = 0,\,\,\,\,
 \mbox{\'{o}}\,\, \mbox{tambi\'{e}n:}\,\,\, \left| {{\begin{array}{*{20}c}
  {x} \hfill & {y} \hfill & {z} \hfill & {1} \hfill  \\
  {a_{1}^{} } \hfill & {a_{2}^{} } \hfill & {a_{3}^{} } \hfill & {1} \hfill
\\
  {b_{1}^{} } \hfill & {b_{2}^{} } \hfill & {b_{3}^{} } \hfill & {1} \hfill
\\
  {c_{1}^{} } \hfill & {c_{2}^{} } \hfill & {c_{3}^{} } \hfill & {1} \hfill
\\
\end{array} }} \right| = 0.\]
Para ver esto obs\'{e}rvese que si se le resta la 2da. fila a la
1ra., la 3da. y la 4ta. filas, el  determinante no cambia.
As\'{\i} se tiene:
\[
\left| {{\begin{array}{*{20}c}
  {x - a_{1}^{} } \hfill & {y - a_{2}^{} } \hfill & {z - a_{3}^{} } \hfill &
{0} \hfill  \\
  {a_{1}^{} } \hfill & {a_{2}^{} } \hfill & {a_{3}^{} } \hfill & {1} \hfill
\\
  {b_{1}^{}  - a_{1}^{} } \hfill & {b_{2}^{}  - a_{2}^{} } \hfill &
{b_{3}^{}  - a_{3}^{} } \hfill & {0} \hfill  \\
  {c_{1}^{}  - a_{1}^{} } \hfill & {c_{2}^{}  - a_{2}^{} } \hfill &
{c_{3}^{}  - a_{3}^{} } \hfill & {0} \hfill  \\
\end{array} }} \right|\,\, = \,\,
\left| {{\begin{array}{*{20}c}
  {x - a_{1}^{} } \hfill & {y - a_{2}^{} } \hfill & {z - a_{3}^{} } \hfill
\\
  {b_{1}^{}  - a_{1}^{} } \hfill & {b_{2}^{}  - a_{2}^{} } \hfill &
{b_{3}^{}  - a_{3}^{} } \hfill  \\
  {c_{1}^{}  - a_{1}^{} } \hfill & {c_{2}^{}  - a_{2}^{} } \hfill &
{c_{3}^{}  - a_{3}^{} } \hfill  \\
\end{array} }} \right| = 0.
\]

%\end{document}

%ch1.tex


\chapter{Basics of Extension and Lifting Problems}
\begin{center}
{\small\em To boldly go where no map has gone before}
\end{center}

\section{Existence problems}
We begin with some metamathematics.
All problems about the existence of maps can be cast into one of the
following two forms, which are in a sense mutually dual.

\noindent
{\bf The Extension Problem}\index{extension problem} \    %%% NB index entry tag
Given an inclusion $ A \stackrel{i}{\hookrightarrow} X $, and a map
$ A \stackrel{f}{\rightarrow} Y $,
does there exist a map $f^{\dagger}:X\to Y$ such that
$f^{\dagger}$ agrees with $f$ on $A$?

Here the appropriate source category for maps should be clear from the
context and, moreover, commutativity through a
candidate $f^{\dagger}$ is precisely
the restriction requirement; that is,
$$f^{\dagger}   :  f^{\dagger}\circ i = f^{\dagger}|_A = f\,. $$
If such an $f^{\dagger}$ exists\footnote{${}^{\dagger}$ suggests striving
for perfection, crusading}, then it is called an {\bf
extension}\index{extension!of a map|bi} of $f$ and is said to {\bf
extend}\index{extend|bi} $f$. In any diagrams, the presence of
a dotted arrow or an arrow carrying a ? indicates a pious hope, in no way
begging the question of its existence. Note that we shall usually
omit $\circ$ from composite maps.

\noindent
{\bf The Lifting Problem}\index{lifting problem} \
Given a pair of maps $E \stackrel{p}{\rightarrow}B$ and $X \stackrel{f}
{\rightarrow} B $,
does there exist a map $f^{\circ} : X \to E$, with
$pf^{\circ} = f  $?


That {\em all\/} existence problems about maps are essentially of one
type or
the other from these two is seen as follows. Evidently, all existence problems
are representable by triangular diagrams\index{triangular diagrams} and it
is easily seen that there are only these six possibilities:
\begin{center}\begin{picture}(300,70)  %augch2 75
\put(5,60){\vector(1,0){30}}
\put(55,60){\vector(1,0){30}}
\put(135,60){\vector(-1,0){30}}
\put(185,60){\vector(-1,0){30}}
\put(235,60){\vector(-1,0){30}}
\put(285,60){\vector(-1,0){30}}
\put(0,55){\vector(0,-1){30}}
\put(50,55){\vector(0,-1){30}}
\put(100,25){\vector(0,1){30}}
\put(150,25){\vector(0,1){30}}
\put(200,55){\vector(0,-1){30}}
\put(250,55){\vector(0,-1){30}}
\put(28,33){\small ?}
\put(78,33){\small ?}
\put(128,33){\small ?}
\put(178,33){\small ?}
\put(228,33){\small ?}
\put(278,33){\small ?}
\put(10,3){\bf 1}
\put(60,3){\bf 2}
\put(110,3){\bf 3}
\put(160,3){\bf 4}
\put(210,3){\bf 5}
\put(260,3){\bf 6}
\put(35,55){\vector(-1,-1){30}}
\put(155,25){\vector(1,1){30}}
\put(135,55){\vector(-1,-1){30}}
\put(55,25){\vector(1,1){30}}
\put(235,55){\vector(-1,-1){30}}
\put(255,25){\vector(1,1){30}}
\end{picture}\end{center}



%\begin{figure}
%\begin{picture}(300,220)(0,0)
%\put(-20,-20){\resizebox{20 cm}{!}{\includegraphics{3dpdf}}}
%\put(260,-10){\resizebox{15 cm}{!}{\includegraphics{contpdf}}}
%\put(220,80){$\beta$}
%\put(400,-10){$N$}
%\put(260,170){$\beta$}
%\put(90,15){$N$}
%\end{picture}
%\caption{{\em The log-gamma family of densities with central mean
%$<N> \, = \frac{1}{2}$ as a surface and as a contour plot. }}
%\label{pdf}
%\end{figure}

\newpage

%\include{3.apendice3}
%\clearemptydoublepage
%Otros apendices
%\bibliographystyle{abbrv}
%\bibliography{myblib}
\begin{thebibliography}{}

\bibitem{A-P}\label{AP}
Álvarez, F.; Peypouquet, J.
\newblock{\em Introducci\'on a la Teor\'ia de Semigrupos}
\newblock Apuntes para la III Escuela de Verano; Universidad de Chile, 2003.

\bibitem{A-H}\label{AH}
Atkinson, K.; Han, W.
\newblock{\em Theoretical Numerical Analysis: a functional analysis framework}
\newblock Texts in applied mathematics; 39.

\bibitem{B-F}\label{BF}
Burden, R.; Faires, J.
\newblock{\em An\'alisis Num\'erico}
\newblock Thomson Learning, Mexico, D.F. 2002.

\bibitem{Br}\label{Br}
Brezis, H.
\newblock{\em Analyse fonctionnelle: the\'orie et applications}
\newblock Masson, Paris, 1983

\bibitem{Ca}\label{Ca}
Carbajal P., Efra\'in
\newblock{\em Una clase de algor\'itmos de alta precisi\'on para un sistema acoplado}.
\newblock PESQUIMAT, Revista de la Unidad de Post Grado - UNMSM, Vol. 3, Nro. 1, Junio 1999, pag. 1-14.

\bibitem{Cod}\label{Cod}
Code\c{c}o Coelho, F.
\newblock{\em Computa\c{c}ão cient\'ifica com Python : Uma introdu\c{c}ão �  programa\c{c}ão para cientistas}
\newblock Petr\'opolis - RJ, 2007.

\bibitem{Cun}\label{Cun}
Cunha, C. A.
\newblock{\em Semigrupos Aplicados a Sistemas Dissipativos em EDP}
\newblock Florian\'opolis, SC: SBMAC, 2007.


\bibitem{Daf}\label{Daf}
Dafermos, C. M.
\newblock{\em On the existence and the asymptotic stability of solution to the equatios of linear thermoelasticity}
\newblock Arch. Rational Mech. Anal., 29 (1968), 241-271.

\bibitem{Hoff}\label{Hoff}
Hoffman, Joe D.
\newblock{\em Numerical methods for engineers and scientists}
\newblock Marcel Dekker, Inc. 2001.

\bibitem{Ji-Rac}\label{Ji-Rac}
Jiang, Song; Racke, Reinhard
\newblock{\em Evolution Equations in Thermoelasticity}.
\newblock Chapman \& Hall/CRC. 2000.

\bibitem{Jo}\label{Jo}
Johnson, A.
\newblock{\em The Fundamental Theorem of Numerical Analysis}.
\newblock 2006.

\bibitem{Kin}\label{Kin}
Kincaid, David; Ward, E.   
\newblock{\em Numerical analysis : mathematics of scientific computing}
\newblock Brooks/Cole Publishing Company. 1990.

\bibitem{Kr}\label{Kr}
Kreyszig, E.
\newblock{\em Introductory functional analysis with applications.}
\newblock John Wiley e Sons, 1978.

\bibitem{Mu}\label{Mu}
Mu\~noz Rivera, J.E.
\newblock{\em Semigrupos e Equa\c{c}\~{o}es Diferenciais Parciais}
\newblock Laborat\'orio Nacional de Computa\c{c}ão Cient\'ifica, Textos de p\'os gradua\c{c}ão.

\bibitem{Mu-1}\label{Mu1}
Mu\~noz Rivera, J.E.
\newblock{\em Energy decay rates in linear thermoelasticity}
\newblock Funkcialaj Ekvacioj, 35(1) (1992), 19-30.

\bibitem{Pa}\label{Pa}
Pazy, A.
\newblock{\em Semigroups of linear operators and applications to partial differential equations}
\newblock Springer-Verlag, New York 1983.

\bibitem{Pe-Mz}\label{Pe-Mz}
Pereira, D.C.; Menzala, G.P.
\newblock{\em Exponential decay of solutions in thermoelasticity}
\newblock Comp. Appl. Math., 8(1989), 192-204.

\bibitem{Ra}\label{Ra}
Racke, R.
\newblock{\em Exponential decay for a class of initial boundary value problems in thermoelasticity}
\newblock Comp. Appl. Math., 12(1993), 67-80.

\bibitem{R-K}\label{RK}
Richtmyer, R.D.; Morton, K.W.
\newblock{\em Difference Methods for Initial Value Problems}
\newblock Interscience Publishers, New York, 1967.

\bibitem{St-Bu}\label{StBu}
Stoer, J.; Bulirsch, R.
\newblock{\em Introduction to numerical analysis}
\newblock Springer-Verlag New York, Inc. 1992.

\bibitem{Yo}\label{Yo}
Yosida, Kôsaku 
\newblock{\em Functional Analysis}.
\newblock Splinger-Verlag, Berlin Heidelberg New York 1980.

\bibitem{Zua}\label{Zua}
Zuazua, Enrique 
\newblock{\em T\'opicos de An\'alisis Num\'erico}.
\newblock Departamento de Matem\'aticas, Universidad Aut\'onoma. Madrid, Espa\~na.

% Brezis -- kreyzig




\end{thebibliography}


\clearemptydoublepage
\end{document}